\begin{textsample}{NEG Dim 2 – vem\_ed\_10 – Score 1.00 – vem\_ed\_10/t043.txt}  \label{ex:f2_neg_034}
Jornada de PCIs: compartilhando experiências “We are SOME of the champions” reuniu relatos apresentados por professores com experiência consolidada em diversos tipos de PCIs.
Camila Kami, professora de inglês, tem experiências com PCIs desde 2019, nas Fatecs Bauru e Garça.
Ressaltou a importância de ferramentas tecnológicas user-friendly, como o Padlet.
Na atividade de ice-breaker, sugeriu um piquenique virtual.
Recomendou reuniões semanais com o professor parceiro, conversas com os alunos sobre questões interculturais e rubricas com critérios de avaliação.
Elizabeth Herrera Colorado, professora de espanhol, relatou um projeto sobre interculturalidade da Fatec Itaquaquecetuba com a UDEM (México). “Paciência, flexibilidade e acompanhamento permanente são nossos aliados”, resumiu.
Ela recomenda ainda entrar em sintonia com os parceiros internacionais e preparar os alunos para as apresentações síncronas.
Ricardo Pompeu comentou o PCI que as Fatecs Americana e Sumaré desenvolveram com Symbiosis sobre satisfação no trabalho, envolvendo 100 alunos (50 de cada país).
Segundo Pompeu, preparação prévia, tanto dos professores quanto dos alunos, foi essencial para o sucesso. “O PCI vai muito além da ementa da disciplina”, frisou.
Neusa Haruka contou sobre o PCI com Tianjin Normal University (China), sobre aprendizado do inglês como língua estrangeira.
O desafio principal é o fuso horário de 11 horas, o que permitia reuniões síncronas apenas nos fins de semana.
Mesmo assim, os alunos participavam.
A atividade intercultural consiste em pesquisas sobre a cultura do outro país.
Esse projeto tem três anos e meio e envolve 8 Fatecs.

% matched lemmas: 
\end{textsample}
